\documentclass[12pt]{article}
\usepackage[utf8]{inputenc}
\usepackage[spanish]{babel}
\usepackage{amsmath}
\usepackage{amsfonts}
\usepackage{amssymb}
\usepackage{geometry}
\usepackage{fancyhdr}
\usepackage{parskip}
\usepackage{float}
\usepackage{hyperref}
\usepackage{tikz}
\usetikzlibrary{arrows.meta,positioning}
\AtBeginEnvironment{tikzpicture}{\shorthandoff{<>}}

\geometry{a4paper, margin=1in}

\pagestyle{fancy}
\fancyhf{}
\lhead{Astroestadística }
\chead{Problem Set 5}
\rhead{Problemas Pares}


\begin{document}

\section*{Problema 2}

\textbf{Elecciones.} Dos partidos políticos G y U han comenzado su campaña electoral en el país. Sea $p$ la proporción de todos los votantes potenciales que favorecen a U sobre G. Considere probar $H_0: p = 0.5$ contra $H_a: p \neq 0.5$ basado en una muestra aleatoria de 100 individuos. Sea el estadístico de prueba $X$ el número de la muestra que favorece a U y $x$ representa el valor observado de $X$.

\begin{enumerate}
    \item[a)] Describa los errores tipo I y tipo II en este contexto.
    \item[b)] Suponga que $x = 37$. ¿Qué valores de $X$ son igual o más contradictorios que éste con $H_0$?
    \item[c)] ¿Cuál es la distribución de probabilidades del estadístico de prueba $X$ cuando $H_0$ es cierto? Use esto para calcular la probabilidad de obtener un valor igual o más contradictorio a $H_0$ cuando este es cierto (P-value) si $x = 37$.
    \item[d)] Si se va a rechazar $H_0$ cuando P-value $\leq 0.056$, calcule la probabilidad de tener un error de tipo II cuando $p$ tiene valores de 0.4, 0.3, 0.6 y 0.7. [Ayuda: ¿A qué desigualdades que involucran a $x$ es equivalente P-value $> 0.056$?]
    \item[e)] Usando el procedimiento de prueba de (d), ¿qué concluiría si 37 de los 100 consultados favorecieran a U?
\end{enumerate}

\section*{Solución}

\subsection*{a): Errores Tipo I y Tipo II}

Debemos definir los errores en el contexto del problema:

\textbf{Error Tipo I:} Ocurre cuando rechazamos la hipótesis nula $H_0$ siendo verdadera. En este contexto electoral:

\begin{quote}
    \textit{``Concluir que la proporción de votantes que favorecen a U es diferente de 0.5 (hay un candidato preferido), cuando en realidad la proporción es exactamente 0.5 (no hay preferencia entre los candidatos).''}
\end{quote}

Matemáticamente, el error tipo I se denota como $\alpha = P(\text{rechazar } H_0 | H_0 \text{ es verdadera})$.

\textbf{Error Tipo II:} Ocurre cuando no rechazamos la hipótesis nula $H_0$ siendo falsa. En este contexto:

\begin{quote}
    \textit{``Concluir que no hay diferencia en la preferencia de votantes (proporción = 0.5), cuando en realidad sí existe una preferencia (la proporción es diferente de 0.5).''}
\end{quote}

Matemáticamente, el error tipo II se denota como $\beta = P(\text{no rechazar } H_0 | H_0 \text{ es falsa})$.

\subsection*{b): Valores Igual o Más Contradictorios}

Dado que $x = 37$ y $n = 100$, bajo $H_0: p = 0.5$, el valor esperado es:
$$E(X) = np = 100(0.5) = 50$$

La desviación observada respecto al valor esperado es:
$$|37 - 50| = 13$$

Para una prueba de dos colas, los valores igual o más contradictorios son aquellos que están al menos tan lejos del valor esperado como la observación. Por lo tanto, buscamos valores de $X$ tales que:
$$|X - 50| \geq 13$$

Esto implica:
$$X - 50 \geq 13 \quad \text{o} \quad X - 50 \leq -13$$
$$X \geq 63 \quad \text{o} \quad X \leq 37$$

\textbf{Respuesta:} Los valores igual o más contradictorios son $X \in \{0, 1, 2, \ldots, 37\} \cup \{63, 64, \ldots, 100\}$.

\begin{figure}[H]
\centering
\begin{tikzpicture}[x=0.12cm,y=0.9cm,>=Stealth]
    % Baseline
    \draw[->] (0,0) -- (101,0) node[right] {$X$};

    % Regions
    \fill[red!18] (0,0) rectangle (37,0.7);
    \fill[blue!10] (37,0) rectangle (63,0.7);
    \fill[red!18] (63,0) rectangle (100,0.7);

    % Ticks and labels
    \foreach \x/\lab in {0/0,20/20,37/37,50/50,63/63,80/80,100/100}
        \draw (\x,0.08) -- (\x,-0.12) node[below] {\lab};

    % Boundary markers
    \draw[very thick,red!70!black] (37,0) -- (37,1.05);
    \draw[very thick,red!70!black] (63,0) -- (63,1.05);
    \draw[dashed] (50,0) -- (50,1.05);

    % Annotations
    \node[red!70!black,align=center] at (18.5,0.35) {$X\leq 37$};
    \node[align=center] at (50,0.35) {$38\leq X\leq 62$};
    \node[red!70!black,align=center] at (81.5,0.35) {$X\geq 63$};
\end{tikzpicture}
\caption{En una prueba bilateral, los valores igual o más contradictorios que $x=37$ son los que quedan al menos a 13 unidades de 50. Por ello, las regiones extremas son $X\leq 37$ y $X\geq 63$.}
\end{figure}

\subsection*{c): Distribución y P-value}

\textbf{Distribución de Probabilidad y Justificación de la Aproximación Normal:}

Cuando $H_0: p = 0.5$ es verdadera, el estadístico $X$ sigue una \textbf{distribución binomial} con parámetros $n = 100$ y $p = 0.5$:
$$X \sim \text{Binomial}(n=100, p=0.5)$$

Verificamos si podemos usar la aproximación normal. Los criterios son:
$$np = 100(0.5) = 50 \geq 10 \quad \text{y} \quad n(1-p) = 100(0.5) = 50 \geq 10$$

Ambas condiciones se satisfacen, por lo que podemos aproximar la distribución binomial con una \textbf{distribución normal}:
$$X \sim N(\mu = np, \sigma^2 = np(1-p))$$
$$X \sim N(\mu = 50, \sigma^2 = 25)$$

Por lo tanto:
$$\sigma = \sqrt{25} = 5$$

\textbf{Cálculo del P-value usando la Aproximación Normal:}

El P-value es la probabilidad de obtener un resultado igual o más extremo que el observado, asumiendo que $H_0$ es verdadera. Para una prueba de dos colas con $x = 37$:

$$\text{P-value} = P(X \leq 37) + P(X \geq 63)$$

Estandarizamos usando el estadístico $Z$:
$$Z = \frac{X - \mu}{\sigma} = \frac{37 - 50}{5} = \frac{-13}{5} = -2.6$$

Para el otro extremo:
$$Z = \frac{63 - 50}{5} = \frac{13}{5} = 2.6$$

Por simetría de la distribución normal estándar:
$$\text{P-value} = P(Z \leq -2.6) + P(Z \geq 2.6) = 2 \cdot P(Z \leq -2.6)$$

Consultando la tabla de la distribución normal estándar:
$$P(Z \leq -2.6) = \Phi(-2.6) \approx 0.00466$$

Por lo tanto:
$$\text{P-value} = 2(0.00466) \approx 0.00932 \approx 0.009$$

\begin{figure}[H]
\centering
\begin{tikzpicture}[x=0.16cm,y=40cm,>=Stealth]
    % Axes
    \draw[->] (20,0) -- (80,0) node[right] {$X$};
    \draw[->] (20,0) -- (20,0.085) node[above] {probabilidad};

    % Ticks
    \foreach \x in {20,30,37,40,50,60,63,70,80}
        \draw (\x,0) -- (\x,-0.0018) node[below] {\x};

    % Stylized bell curve centered at 50
    \draw[thick,blue!70!black,smooth,samples=100,domain=20:80]
        plot (\x,{0.08*exp(-((\x-50)*(\x-50))/120)});

    % Shade tails
    \fill[red!20,domain=20:37,smooth,samples=60]
        (20,0) -- plot (\x,{0.08*exp(-((\x-50)*(\x-50))/120)}) -- (37,0) -- cycle;
    \fill[red!20,domain=63:80,smooth,samples=60]
        (63,0) -- plot (\x,{0.08*exp(-((\x-50)*(\x-50))/120)}) -- (80,0) -- cycle;

    % Observed/critical markers
    \draw[dashed,red!70!black] (37,0) -- (37,{0.08*exp(-((37-50)*(37-50))/120)});
    \draw[dashed,red!70!black] (63,0) -- (63,{0.08*exp(-((63-50)*(63-50))/120)});
    \draw[dashed] (50,0) -- (50,0.08);

    \node[red!70!black,align=center] at (29,0.028) {$P(X\leq 37)$};
    \node[red!70!black,align=center] at (71,0.028) {$P(X\geq 63)$};
    \node[below right] at (50,0.08) {$\mu=50$};
\end{tikzpicture}
\caption{Representación del P-value bilateral bajo $H_0$: el área en ambas colas corresponde a $P(X\leq 37)+P(X\geq 63)$.}
\end{figure}


\subsection*{d): Probabilidad de Error Tipo II}

Nota: Aunque el P-value para $x=37$ es aproximadamente 0.009 (del inciso c), 
ahora consideramos un nivel de significancia diferente de $\alpha = 0.056$ para este problema.

La probabilidad de error tipo II, $\beta(p)$, depende del valor verdadero de $p$.

\textbf{Determinación de la región de no rechazo usando Aproximación Normal:}

Se rechaza $H_0$ si P-value $\leq 0.056$. Equivalentemente, no se rechaza si P-value $> 0.056$.

Para una prueba de dos colas con P-value $= 0.056$, cada cola tiene probabilidad $0.028$. Buscamos el valor crítico $z_{\alpha/2}$ tal que:
$$P(Z \leq -z_{0.028}) = 0.028$$

Consultando la tabla normal estándar, encontramos que $z_{0.028} \approx 1.90$ (interpolando entre $z_{0.025} = 1.96$ y $z_{0.03} = 1.88$).

Bajo $H_0$ con $p = 0.5$, $\mu = 50$ y $\sigma = 5$, los valores críticos en términos de $X$ son:
$$X = \mu \pm z_{\alpha/2} \cdot \sigma = 50 \pm 1.90(5) = 50 \pm 9.5$$

Por lo tanto:
$$X \leq 40.5 \quad \text{o} \quad X \geq 59.5$$

Redondeando a números enteros:
$$X \leq 40 \quad \text{o} \quad X \geq 60$$

\textbf{Región de no rechazo:} $41 \leq X \leq 59$ (o equivalentemente, rechazamos si $X \leq 40$ o $X \geq 60$)

\textbf{Cálculo de $\beta(p)$ usando Aproximación Normal:}

Para cada valor de $p$, aproximamos $X \sim N(\mu = np, \sigma = \sqrt{np(1-p)})$ y calculamos:
$$\beta(p) = P(41 \leq X \leq 59 | p)$$

\begin{enumerate}
    \item \textbf{Cuando $p = 0.4$:}
    $$X \sim N(\mu = 40, \sigma = \sqrt{100(0.4)(0.6)} = 4.90)$$
    
    Estandarizamos los límites:
    $$Z_{\text{inf}} = \frac{41 - 40}{4.90} = \frac{1}{4.90} \approx 0.204$$
    $$Z_{\text{sup}} = \frac{59 - 40}{4.90} = \frac{19}{4.90} \approx 3.878$$
    
    $$\beta(0.4) = P(0.204 \leq Z \leq 3.878) = \Phi(3.878) - \Phi(0.204)$$
    $$\beta(0.4) \approx 0.9999 - 0.5808 \approx 0.419$$
    
    \item \textbf{Cuando $p = 0.3$:}
    $$X \sim N(\mu = 30, \sigma = \sqrt{100(0.3)(0.7)} = 4.58)$$
    
    Estandarizamos:
    $$Z_{\text{inf}} = \frac{41 - 30}{4.58} = \frac{11}{4.58} \approx 2.402$$
    $$Z_{\text{sup}} = \frac{59 - 30}{4.58} = \frac{29}{4.58} \approx 6.333$$
    
    $$\beta(0.3) = P(2.402 \leq Z \leq 6.333) = \Phi(6.333) - \Phi(2.402)$$
    $$\beta(0.3) \approx 0.9999 - 0.9918 \approx 0.008$$
    
    \item \textbf{Cuando $p = 0.6$:}
    Por simetría con respecto a $p = 0.5$ (es decir, $0.6 = 0.5 + 0.1$ y $0.4 = 0.5 - 0.1$):
    $$\beta(0.6) = \beta(0.4) \approx 0.419$$
    
    \item \textbf{Cuando $p = 0.7$:}
    Por simetría con respecto a $p = 0.5$:
    $$\beta(0.7) = \beta(0.3) \approx 0.008$$
\end{enumerate}


\textbf{Tabla Resumen de Potencia del Contraste (Aproximación Normal):}

\begin{center}
\begin{tabular}{|c|c|c|}
\hline
\textbf{Valor de } $p$ & \textbf{} $\beta(p)$ & \textbf{Potencia } $1-\beta(p)$ \\
\hline
$0.3$ & $0.008$ & $0.992$ \\
$0.4$ & $0.419$ & $0.581$ \\
$0.6$ & $0.419$ & $0.581$ \\
$0.7$ & $0.008$ & $0.992$ \\
\hline
\end{tabular}
\end{center}

Esta tabla muestra que el contraste tiene excelente potencia para detectar desviaciones grandes de $p = 0.5$ (como $p = 0.3$ o $p = 0.7$), pero potencia moderada para detectar desviaciones menores (como $p = 0.4$ o $p = 0.6$).

\begin{figure}[H]
\centering
\begin{tikzpicture}[x=9cm,y=6cm,>=Stealth]
    % Axes
    \draw[->] (0.26,0) -- (0.74,0) node[right] {$p$};
    \draw[->] (0.26,0) -- (0.26,1.08) node[above] {$\beta(p)$};

    % Ticks
    \foreach \x/\lab in {0.3/0.3,0.4/0.4,0.5/0.5,0.6/0.6,0.7/0.7}
        \draw (\x,0) -- (\x,-0.02) node[below] {\lab};
    \foreach \y/\lab in {0.0/0,0.2/0.2,0.4/0.4,0.6/0.6,0.8/0.8,1.0/1}
        \draw (0.26,\y) -- (0.252,\y) node[left] {\lab};

    % Symmetry line
    \draw[dashed] (0.5,0) -- (0.5,1.05);
    \node[above right] at (0.5,0.95) {simetría en $p=0.5$};

    % Polyline and points
    \draw[thick,blue!70!black] (0.3,0.008) -- (0.4,0.419) -- (0.6,0.419) -- (0.7,0.008);
    \foreach \x/\y in {0.3/0.008,0.4/0.419,0.6/0.419,0.7/0.008}
        \filldraw[blue!70!black] (\x,\y) circle (1.2pt);

    % Labels
    \node[above left,fill=white,inner sep=1pt] at (0.4,0.419) {$(0.4,0.419)$};
    \node[above right,fill=white,inner sep=1pt] at (0.6,0.419) {$(0.6,0.419)$};
    \node[above left,fill=white,inner sep=1pt] at (0.3,0.008) {$(0.3,0.008)$};
    \node[above right,fill=white,inner sep=1pt] at (0.7,0.008) {$(0.7,0.008)$};
\end{tikzpicture}
\caption{Probabilidad de error tipo II para varios valores de $p$ usando la regla de decisión del inciso (d). El error tipo II es mayor cuando $p$ está más cerca de $0.5$.}
\end{figure}



\subsection*{e): Conclusión con $x = 37$}

Del inciso (d), sabemos que:

$\bullet$ Región de rechazo: $X \leq 40$ o $X \geq 60$

$\bullet$ Región de no rechazo: $41 \leq X \leq 59$

Dado que $x = 37$:
$$37 \leq 40 \text{ (está en la región de rechazo)}$$

\begin{figure}[H]
\centering
\begin{tikzpicture}[x=0.14cm,y=0.9cm,>=Stealth]
    \draw[->] (0,0) -- (101,0) node[right] {$X$};
    \foreach \x/\lab in {0/0,20/20,37/37,40/40,50/50,60/60,80/80,100/100}
        \draw (\x,0) -- (\x,-0.18) node[below] {\lab};

    \fill[red!18] (0,0) rectangle (40,0.9);
    \fill[green!18] (40,0) rectangle (60,0.9);
    \fill[red!18] (60,0) rectangle (100,0.9);

    \draw[thick] (40,0) -- (40,0.9);
    \draw[thick] (60,0) -- (60,0.9);
    \draw[very thick,blue!70!black] (37,0) -- (37,1.2);

    \node at (20,0.45) {rechazo};
    \node at (50,0.45) {no rechazo};
    \node at (80,0.45) {rechazo};
    \node[blue!70!black,above] at (37,1.2) {$x=37$};
\end{tikzpicture}
\caption{Región de decisión del inciso (d): como $x=37$ cae en la región de rechazo, se rechaza $H_0$.}
\end{figure}


Por lo tanto, \textbf{rechazamos $H_0$}. El valor P es aproximadamente 0.009, 
que es menor que 0.056, lo que proporciona evidencia fuerte contra $H_0$.

\textbf{Interpretación:} Con 37 de 100 consultados favoreciendo a U, existe evidencia estadística suficiente para concluir que la proporción de votantes que favorecen a U es diferente de 0.5. Específicamente, como $37 < 50$, los datos sugieren que el partido G tiene mayor preferencia que el partido U en la población de votantes.

\section*{Problema 4}

\textbf{Más propina.} Los camareros de los restaurantes han empleado diversas estrategias para aumentar las propinas. Un artículo publicado el 5 de septiembre de 2005 (New Yorker) reportó que ``en un estudio, una camarera recibió un 50 \% más de propinas cuando se presentó por su nombre que cuando no lo hizo''. Considere los siguientes datos (ficticios) sobre la cantidad de propina como porcentaje de la factura:

\begin{itemize}
    \item \textbf{Presentándose:} $m = 50$, $\bar{x} = 22.63$, $s_1 = 7.82$
    \item \textbf{Sin presentarse:} $n = 50$, $\bar{y} = 14.15$, $s_2 = 6.10$
\end{itemize}

¿Estos datos sugieren que presentarse por su nombre aumenta las propinas en promedio en más del 50\%? Establezca y pruebe las hipótesis relevantes. [Sugerencia: considere el parámetro $\theta = \mu_1 - 1.5\mu_2$.]

\section*{Solución}

Para resolver este problema, utilizaremos la teoría de inferencia estadística para dos muestras independientes con tamaños grandes, tal como se detalla en el \textbf{Capítulo 9 (Inferencias basadas en dos muestras)} del texto de Jay L. Devore. Específicamente, adaptaremos la prueba de diferencia de medias para una combinación lineal de medias.

\subsection*{1. Planteamiento de las Hipótesis}

El problema nos pide evaluar si la propina promedio cuando el camarero se presenta ($\mu_1$) es más del 50\% mayor que la propina promedio cuando no se presenta ($\mu_2$). Matemáticamente, esto se expresa como $\mu_1 > 1.5\mu_2$.

Siguiendo la sugerencia, definimos el parámetro de interés $\theta$ como la diferencia entre la media del grupo 1 y 1.5 veces la media del grupo 2:
$$ \theta = \mu_1 - 1.5\mu_2 $$

Las hipótesis nula y alternativa son:
\begin{align*}
    H_0: \theta &= 0 \quad (\text{es decir, } \mu_1 = 1.5\mu_2) \\
    H_a: \theta &> 0 \quad (\text{es decir, } \mu_1 > 1.5\mu_2)
\end{align*}
Esta es una prueba de \textbf{cola superior}.

\subsection*{2. Estadístico de Prueba}

Dado que los tamaños de muestra son grandes ($m=50$ y $n=50$), podemos utilizar la aproximación normal (estadístico $Z$) basándonos en el Teorema del Límite Central, tal como se discute en la \textbf{Sección 9.1} de Devore sobre inferencias de muestras grandes.

Necesitamos un estimador para $\theta$. El estimador insesgado natural es:
$$ \hat{\theta} = \bar{X} - 1.5\bar{Y} $$

Para calcular el estadístico de prueba estandarizado, necesitamos la varianza de este estimador. Utilizando las propiedades de la varianza para combinaciones lineales de variables aleatorias independientes (donde $V(aX + bY) = a^2V(X) + b^2V(Y)$):
$$ V(\hat{\theta}) = V(\bar{X} - 1.5\bar{Y}) = 1^2 \cdot V(\bar{X}) + (-1.5)^2 \cdot V(\bar{Y}) $$

Sustituyendo las varianzas muestrales estimadas ($S^2/n$):
$$ \hat{V}(\hat{\theta}) = \frac{s_1^2}{m} + 2.25 \frac{s_2^2}{n} $$

El estadístico de prueba $Z$ bajo $H_0$ (donde $\theta = 0$) es:
$$ Z = \frac{(\bar{x} - 1.5\bar{y}) - 0}{\sqrt{\frac{s_1^2}{m} + 2.25 \frac{s_2^2}{n}}} $$

\subsection*{3. Cálculos Numéricos}

Sustituimos los datos proporcionados:
\begin{itemize}
    \item Grupo 1 (Presentándose): $m = 50, \bar{x} = 22.63, s_1 = 7.82$
    \item Grupo 2 (Sin presentarse): $n = 50, \bar{y} = 14.15, s_2 = 6.10$
\end{itemize}

\textbf{Paso A: Calcular el numerador (diferencia observada)}
$$ \hat{\theta} = \bar{x} - 1.5\bar{y} = 22.63 - 1.5(14.15) = 22.63 - 21.225 = 1.405 $$

\textbf{Paso B: Calcular el error estándar de } $\hat{\theta}$

Primero, calculamos la varianza estimada usando:
$$ \hat{V}(\hat{\theta}) = \frac{s_1^2}{m} + (1.5)^2 \frac{s_2^2}{n} = \frac{s_1^2}{m} + 2.25 \frac{s_2^2}{n} $$

Componente 1 (propinas con presentación):
$$ \frac{s_1^2}{m} = \frac{7.82^2}{50} = \frac{61.1524}{50} = 1.2231 $$

Componente 2 (propinas sin presentación):
$$ 2.25 \times \frac{s_2^2}{n} = 2.25 \times \frac{6.10^2}{50} = 2.25 \times \frac{37.21}{50} = 2.25 \times 0.7442 = 1.6745 $$

Varianza total estimada:
$$ \hat{V}(\hat{\theta}) = 1.2231 + 1.6745 = 2.8976 $$

Error Estándar:
$$ SE(\hat{\theta}) = \sqrt{2.8976} = 1.7023 $$

\textbf{Paso C: Calcular el estadístico Z}
$$ Z = \frac{\hat{\theta}}{SE(\hat{\theta})} = \frac{1.405}{1.7023} = 0.8251 $$


\subsection*{4. Valor P y Decisión}

El valor P es la probabilidad de obtener un estadístico de prueba tan extremo o más extremo que el observado, asumiendo que $H_0$ es verdadera. Al ser una prueba de cola superior:
$$ P\text{-value} = P(Z > 0.8251) = 1 - \Phi(0.8251) $$

Consultando la \textbf{Tabla A.3 (Áreas de la Curva normal estándar)} del apéndice de Devore:
\begin{itemize}
    \item $\Phi(0.82) \approx 0.7939$
    \item $\Phi(0.83) \approx 0.7967$
\end{itemize}

Interpolando para $0.8251$, obtenemos $\Phi(0.8251) \approx 0.7952$.

$$ P\text{-value} \approx 1 - 0.7952 = 0.2048 $$

Comparamos con el nivel de significancia estándar de $\alpha = 0.05$:
$$ 0.2048 > 0.05 \quad \Rightarrow \quad \text{No rechazamos } H_0 $$

\begin{figure}[H]
\centering
\begin{tikzpicture}[x=1.8cm,y=5cm,>=Stealth]
    % Axes
    \draw[->] (-3.2,0) -- (3.4,0) node[right] {$z$};
    \draw[->] (-3.2,0) -- (-3.2,0.46) node[above] {densidad};

    % Bell curve
    \draw[thick,blue!70!black,smooth,samples=120,domain=-3.2:3.2]
        plot (\x,{0.4*exp(-\x*\x/2)});

    % Shade upper tail
    \fill[red!20,domain=0.825:3.2,smooth,samples=80]
        (0.825,0) -- plot (\x,{0.4*exp(-\x*\x/2)}) -- (3.2,0) -- cycle;

    % Markers
    \draw[dashed] (0,0) -- (0,0.4);
    \draw[dashed,red!70!black] (0.825,0) -- (0.825,{0.4*exp(-0.825*0.825/2)});

    % Ticks
    \foreach \x/\lab in {-3/-3,-2/-2,-1/-1,0/0,2/2,3/3}
        \draw (\x,0) -- (\x,-0.012) node[below] {\lab};
    \draw (0.825,0) -- (0.825,-0.012) node[below left] {0.825};
    \draw (1,0) -- (1,-0.012) node[below right] {1};

    % Labels
    \node[red!70!black] at (1.95,0.11) {$P(Z>0.8251)$};
\end{tikzpicture}
\caption{Representación gráfica del valor P para la prueba de cola superior del problema 4. El área sombreada corresponde a $P(Z>0.8251) \approx 0.205$.}
\end{figure}

Por lo tanto, \textbf{no rechazamos la hipótesis nula} al nivel de significancia del 5\%.

\subsection*{5. Conclusión Final}

Con un valor P de aproximadamente 0.205, no existe evidencia estadística suficiente para rechazar la hipótesis nula $H_0: \mu_1 = 1.5\mu_2$ al nivel de significancia del 5\%. Los datos \textbf{no sugieren} que presentarse por nombre aumente las propinas en promedio en más del 50\%.

\textbf{Observación:} Aunque la propina promedio cuando el camarero se presentó fue de $22.63\%$, comparada con $1.5 \times 14.15\% = 21.225\%$ (1.5 veces la propina sin presentarse), esta diferencia de tan solo $1.405$ puntos porcentuales no es estadísticamente significativa. La variabilidad muestral es demasiado grande como para concluir que hay un aumento real del 50\% en las propinas. Se requeriría un tamaño de muestra mayor o un efecto más pronunciado para detectar estadísticamente esta diferencia.

\section*{Problema 6}

\textbf{Cúmulos jóvenes.} Suponga que $\mu_1$ y $\mu_2$ son las medias reales de las edades de dos cúmulos estelares jóvenes. Por años se ha considerado que el primer cúmulo es diez años más joven que el segundo. Sin embargo, un estudio reciente ha encontrado datos que pueden sugerir que el primero es aún más joven. Use un nivel de significancia de 0.01 para probar esto último, si se tienen los siguientes datos: $m = 25$, $\bar{x} = 115.7$, $s_1 = 5.03$, para el primer cúmulo y $n = 31$, $\bar{y} = 129.3$, y $s_2 = 5.38$ para el segundo.

\section*{Solución}

Para resolver este problema, utilizaremos los procedimientos de inferencia para dos muestras basados en la distribución $t$, específicamente la prueba $t$ para dos muestras con varianzas desiguales (aproximación de Welch-Satterthwaite), tal como se describe en la \textbf{Sección 9.2} del texto de referencia (Devore).

\subsection*{1. Definición de Parámetros e Hipótesis}

Sean:
\begin{itemize}
    \item $\mu_1$: la edad media real del primer cúmulo estelar.
    \item $\mu_2$: la edad media real del segundo cúmulo estelar.
\end{itemize}

El problema establece que históricamente se consideraba que el primer cúmulo era 10 años más joven que el segundo, es decir, $\mu_2 - \mu_1 = 10$. El estudio reciente sugiere que el primer cúmulo es \textit{aún más joven}, lo que implica que la diferencia de edad es mayor a 10 años ($\mu_2 - \mu_1 > 10$).

Por lo tanto, planteamos las hipótesis de la siguiente manera:
\begin{align*}
    H_0 &: \mu_2 - \mu_1 = 10 \\
    H_a &: \mu_2 - \mu_1 > 10
\end{align*}
Esta es una prueba de cola superior.

\begin{figure}[H]
\centering
\begin{tikzpicture}[x=1.2cm,y=3.6cm,>=Stealth]
    % Axes
    \draw[->] (-4,0) -- (4.2,0) node[right] {$d=\mu_2-\mu_1$};
    \draw[->] (0,0) -- (0,1.15) node[above] {plausibilidad};

    % Reference line at null value
    \draw[dashed] (0,0) -- (0,1.02);
    \node[below] at (0,0) {$10$};

    % Alternative direction
    \draw[very thick,blue!70!black,->] (0.35,0.62) -- (3.55,0.62);
    \node[blue!70!black,align=center] at (2.2,0.78) {$H_a:\,\mu_2-\mu_1>10$\\cola superior};

    % Null region marker
    \node[align=center] at (-1.8,0.42) {$H_0:\,\mu_2-\mu_1=10$};
\end{tikzpicture}
\caption{Esquema de las hipótesis: la afirmación histórica fija la diferencia en $10$ años, mientras que la evidencia buscada apunta a valores mayores que $10$.}
\end{figure}

\subsection*{2. Estadístico de Prueba}

Dado que las desviaciones estándar muestrales ($s_1 = 5.03$ y $s_2 = 5.38$) son similares pero no idénticas, y los tamaños de muestra son diferentes ($m=25, n=31$), utilizaremos el estadístico $t$ para dos muestras con varianzas no agrupadas. El estadístico de prueba se define como:

$$ t = \frac{(\bar{y} - \bar{x}) - (\mu_2 - \mu_1)_0}{\sqrt{\frac{s_1^2}{m} + \frac{s_2^2}{n}}} $$

Donde $(\mu_2 - \mu_1)_0 = 10$ es el valor de la diferencia bajo la hipótesis nula.

\subsection*{3. Cálculos}

Sustituimos los valores proporcionados en el enunciado:
\begin{itemize}
    \item Muestra 1 (Primer cúmulo): $m = 25, \bar{x} = 115.7, s_1 = 5.03$
    \item Muestra 2 (Segundo cúmulo): $n = 31, \bar{y} = 129.3, s_2 = 5.38$
\end{itemize}

Primero calculamos el error estándar de la diferencia:
$$ SE = \sqrt{\frac{s_1^2}{m} + \frac{s_2^2}{n}} = \sqrt{\frac{(5.03)^2}{25} + \frac{(5.38)^2}{31}} $$
$$ SE = \sqrt{\frac{25.3009}{25} + \frac{28.9444}{31}} = \sqrt{1.0120 + 0.9337} = \sqrt{1.9457} \approx 1.3949 $$

Ahora calculamos el valor del estadístico $t$:
$$ t = \frac{(129.3 - 115.7) - 10}{1.3949} = \frac{13.6 - 10}{1.3949} = \frac{3.6}{1.3949} \approx 2.58 $$

\subsection*{4. Grados de Libertad y Valor Crítico}

Utilizamos la aproximación de Satterthwaite para los grados de libertad ($\nu$):
$$ \nu = \frac{\left( \frac{s_1^2}{m} + \frac{s_2^2}{n} \right)^2}{\frac{(s_1^2/m)^2}{m-1} + \frac{(s_2^2/n)^2}{n-1}} $$

Sustituyendo los valores intermedios calculados ($A = 1.0120, B = 0.9337$):
$$ \nu = \frac{(1.9457)^2}{\frac{(1.0120)^2}{24} + \frac{(0.9337)^2}{30}} = \frac{3.7857}{\frac{1.0241}{24} + \frac{0.8718}{30}} $$
$$ \nu = \frac{3.7857}{0.0427 + 0.0291} = \frac{3.7857}{0.0718} \approx 52.7 $$

Redondeamos a $\nu = 53$ grados de libertad. (Nota: La aproximación de Satterthwaite puede redondearse al entero más cercano; tomamos 53.)

Para un nivel de significancia $\alpha = 0.01$ y una prueba de cola superior con 53 grados de libertad, consultamos la tabla de distribución $t$ de student.

El valor crítico $t_{0.01, 53}$ se encuentra interpolando entre valores tabulados:
\begin{itemize}
    \item $t_{0.01, 50} \approx 2.403$
    \item $t_{0.01, 60} \approx 2.390$
\end{itemize}

Para $\nu = 53$ grados de libertad, interpolamos: $t_{0.01, 53} \approx 2.398 \approx 2.40$.

Valor crítico: $t_{crit} = 2.40$.

\begin{figure}[H]
\centering
\begin{tikzpicture}[x=1.55cm,y=4.2cm,>=Stealth]
    % Axes
    \draw[->] (-3.4,0) -- (4.2,0) node[right] {$t$};
    \draw[->] (-3.4,0) -- (-3.4,0.46) node[above] {densidad};

    % t-like bell curve
    \draw[thick,blue!70!black,smooth,samples=140,domain=-3.4:4]
        plot (\x,{0.4/(1+0.22*\x*\x)^2});

    % Rejection region
    \fill[red!20,domain=2.4:4,smooth,samples=80]
        (2.4,0) -- plot (\x,{0.4/(1+0.22*\x*\x)^2}) -- (4,0) -- cycle;

    % Markers
    \draw[dashed] (0,0) -- (0,0.4);
    \draw[dashed,red!70!black] (2.4,0) -- (2.4,{0.4/(1+0.22*2.4*2.4)^2});
    \draw[dashed,orange!80!black] (2.58,0) -- (2.58,{0.4/(1+0.22*2.58*2.58)^2});

    % Ticks
    \foreach \x/\lab in {-3/-3,-2/-2,-1/-1,0/0,1/1,2/2,3/3,4/4}
        \draw (\x,0) -- (\x,-0.012) node[below] {\lab};
    \draw (2.4,0) -- (2.4,-0.012) node[below=8pt] {$2.40$};
    \draw (2.58,0) -- (2.58,-0.012) node[below right=12pt] {$2.58$};

    % Labels
    \node[red!70!black,align=center] at (3.16,0.11) {región de\\rechazo};
    \node[orange!80!black,align=center] at (2.05,0.31) {$t_{calc}$};
\end{tikzpicture}
\caption{Prueba de cola superior para el problema 6. La región sombreada comienza en el valor crítico $t_{crit}\approx 2.40$, y el estadístico observado $t_{calc}\approx 2.58$ cae dentro de ella.}
\end{figure}

\subsection*{5. Decisión y Conclusión}

Comparamos el estadístico calculado con el valor crítico:
$$ t_{calc} = 2.58 > t_{crit} \approx 2.40 $$

El valor $P$ asociado a $t=2.58$ con 53 grados de libertad es aproximadamente $0.0065$ (consultar tabla $t$ o software estadístico), el cual es menor que $\alpha = 0.01$.

\textbf{Conclusión:} Dado que el estadístico de prueba cae en la región de rechazo:
$$t_{calc} = 2.58 > t_{crit} = 2.40 \quad \text{y} \quad P\text{-value} = 0.0065 < \alpha = 0.01$$

\textbf{Rechazamos la hipótesis nula $H_0$}. Existe evidencia estadística suficiente al nivel de significancia de 0.01 para concluir que el primer cúmulo estelar es significativamente más de 10 años más joven que el segundo cúmulo. En otras palabras, la diferencia de edad excede los 10 años inicialmente establecidos.


\end{document}